\section{Plotinus: The Peak of Pagan Wisdom}

\begin{quotex}
This originally appeared in \emph{Introduction to Magic, Volume 3}, by the \textbf{Gruppo di Ur}. It was published under the name ``Plotino'' and I haven't been able to determine the identity.
\end{quotex}

\begin{quotex}
It is up to the gods to come to me, not for me to go to them.
\end{quotex}


The whole spirit of the ``solar'' way is contained in \textbf{Plotinus}' bold response to Amelio, who invited him to approach the gods with the prescribed rites. The surpassing of the religious attitude, the transcendent dignity of man in possession of Wisdom; and his superiority not only in relation to the natural world, but also in relation to the divine, are asserted.

In particular, it is about an inner meaning, which is essential for practice, because if you do not understand it, if you do not achieve it, it is useless to talk of ``magic'', it is useless to start with abstruse ``ceremonials'', and useless to devote oneself to ascetic disciplines or yogic ``exercises''.

It is necessary to create a quality in oneself, through which the supersensible powers (the gods) are forced to be like females attracted to the male. This quality can be summed up in a word that means nothing and means everything: BEING.

BE, ENDURE, become a CENTER. Through ``ascesis'', through ``purification'', through what Plotinus himself will now make explicit. You have heard of ``solar way''. This is its secret. Separated from those with disordered need, yearning soul, and confused look --- more `non-being' than `being' --- they are attracted to the invisible worlds.

\begin{quotex}
It is necessary to make yourself one self like the gods: not like the leading citizens \emph{uomini da bene} .

Not to be exonerated from sin, but to be a God --- is the goal.
\end{quotex}


These maxims bluntly separate the way of the initiate from the way of men. The ``virtue'' of men, in the final analysis, is something indifferent: ``the image of an image,'' says Plotinus. ``Morality'' has nothing to do with initiation. Initiation is a radical transformation of one state of existence to another state of existence. A ``God'' is not a ``moral value'': it is another being. A good man does not cease to be ``man'' by being ``good.'' At any time and place where the meaning of ``initiation'' is understood, the idea was always the same. Thus in hermetism:

\begin{quotex}
Our work is the conversion and a change of a being into another being, of one thing into another thing, weakness into strength ... corporeality into spirituality.

Even evil men can take water from rivers. Whoever gives simply ignores that which gives.

How does man stand in relation to the whole? As a part? No. As a whole that belongs to himself.

Minus ``one'', are those minus ``being'': more, those who are more.

He is himself, every being, belonging to himself; and to belong to oneself, is to center oneself. The One possesses himself, and has all the greatness and beauty. Behold: he does not flow and does not avoid himself indefinitely. The whole thing is now gathered together in its unity.
\end{quotex}


The first element that constitutes the condition of ``being'' is unity.

\textbf{Unify Yourself --- Be One.}


That bundle of energy, that people of beings, sensations, tendencies that you are, enfold it under a single law, under a single will, under a single thought.

\textbf{Organize Yourself}


Bend your ``soul'', use it in every sense, take it to every crossroad as long as it is inert, incapable of proper motion, dead to every irrationality of instinct. Like a perfectly tamed horse, when driven to the right, it goes right, when driven to the left, it goes left, when braked, it stops, when incited, it flings itself --- so also your soul is for you: one thing you keep wholly in your fist. Without constraints, you will be One: being one, you \emph{are} -- and you belong to yourself. Belonging to yourself, greatness belongs to you.

The ancient classical Aryan wisdom distinguished two symbolic regions: the lower one of the things that ``escape'', the higher, of the ``things that are.'' The things that are powerless to come to the realization and perfect possession of their nature flow and ``escape''. The others are those that have transcended that life that is mixed with death, and that is a flowing and continuous longings. Their ``immobility'', and the same ancient astronomical designation of their ``place'', are symbols. A spiritual state is designated. Being one, no longer dispersed, follows it.

\begin{quotex}
What is the Good for such a man? He is himself his own good. The life that he possesses is perfect. He possesses the good, insofar as he is not in search of another.

To remove whatever is other in respect to his own being, is to purify oneself.

In simple relationship with you without hindrance in your pure unity, without the thing that it is mixed interiorly with this purity, being you only in pure light ... you became a vision.

Although being here, you are ascended.

You no longer need guidance.

Fix your gaze. You'll see.
\end{quotex}


With singular brevity, what is expressed here, in a transcendent sense, is to be called ``good'': the absence of anything that, penetrating inn itself, might take him outside of himself in desire or impulse. Plotinus is careful to clarify the spiritual significance of such a concept by saying that the superior man can even ``look for other things insofar as they are indispensable, not to him, but to those things close to him: to the body to which he is joined, to the life of the body that is not his life. Knowing what the body needs, gives it to him, but these things have no bearing on his life. ''

``Evil'' is the sense of need in the spirit: that of every life that, not knowing how to stand up in himself, loses heart here and there, yearning, looking to complete itself with the achievement of one thing or the other. As long as there is this ``need'', as long as there is this internal and radical insufficiency, the Good is not there. It is nothing nameable: it is an experience that only an act of the spirit on the spirit can produce: separating it from the idea of ?? every other, rejoining it with himself. A state of certainty and fullness then arises, when given, one no longer asks for anything, all talk, all speculation, all agitation, is useless, while one does not know what more can produce a change in the deep mind. Plotinus rightly says that the man, who possesses the whole of his life, possesses perpetuity: being only ``I'', nothing can be added to him, neither in the past nor in the present nor in the future.

\begin{quotex}
The state of being, is to be present.

Every being is in action and in act.

Pleasure is the act of life.

The souls in this universe can also be happy. If they are not -- they accuse themselves, not the universe. They have surrendered in this struggle, where the reward crowns virtue.
\end{quotex}


Plotinus again clarifies the meaning of ``being'': to be is to be present, to be in act. He speaks of ``that intellectual nature without sleep'', a strictly traditional expression. We know the term: ``Awakened'', the ``Always awake'' and the symbolism of the ``sleep'', which furthermore can be more than symbolism in reference to the continuity of a ``being present'' and that does not undergo alteration even in the change of state that usually corresponds to sleep.

To be, therefore, is to be awake. The experience of the whole concentrated in an intellectual clarity, in the simplicity of an act --- it is the experience of `\,''being''. To abandon oneself, to fail --- this is the secret of non-being. The fatigue of the interior unity that loosens and skids, the inner energy that ceases to dominate every part of it as almost, through a disintegration, a multiplicity of tendencies, instincts, and irrational sensations arise -- this is the self-deterioration of the spirit that is manifested in more and more oblique and stunned nature until it reaches the limit form of swoon that is expressed in matter. It is a misunderstanding, says Plotinus, to say that matter is: the being of matter is non-being. Its indefinite divisibility indicates precisely the ``fall'' of the unity that it represents; its inertia, whence it is heavy, strong and blunt, is the same that is characteristic of those who, while failing, cannot hold and deteriorates. That the ``truth'' characteristic of physical knowledge is different does not matter. Corporeal being is the non-being of the spiritual.

Like the actual state of culmination, `\,''being'' is one with the ``good.'' So ``matter'' and ``evil'', in turn, identify with each other; and there is no other ``evil'' except for matter. Here it is naturally necessary to detach oneself, to give up all human preconceptions. The ``evil'' of men has no place in reality and therefore in a metaphysical vision which is always a vision in accordance with reality.

Metaphysically, `good' and `evil' do not exist, but rather what is real and what is not --- and the degree of ``reality'' (meant in the spiritual sense previously explained by ``being'') measures the degree of ``virtue''. In the dry and virile view of ancient classical-Aryan man, only the state of ``privation'' of `\,''being'' was ``bad'': the fatigue, the neglect, the sleep of the interior strength that sets ``matter'' at the limit, as has been said. Neither ``evil'' nor ``matter'' are therefore principles in themselves: they were derived through `degradation' and `dissolution'. Plotinus expresses himself exactly in these terms:

\begin{quotex}
By falling short of the Good, darkness is seen and in the darkness we live. And the evil for the soul is this falling short, the producer of darkness. This is the first evil. Darkness is something that proceeds from it. And the nature of evil is not in matter, but prior to matter in the cessation of the act that gave rise to matter.
\end{quotex}

Plotinus adds: pleasure is the act of life. It is the same view already established by another great genius of the ancient world, Aristotle, who had taught that every activity, insofar as it is perfect, is happy. Happiness and pleasure are like that in the form of purity and freedom: those who arise from the act that is fulfilled and that, once fulfilled, realizes the One, Being, the Good -- not those passive and promiscuous men grabbed into the middle of the turbid self-satisfaction of cravings, thirsts, instincts. Again we are led to the non-human point of view of ``reality''. The one who does not know the irrationality of sentiments. The degree of ``being'' is the secret and measurement of happiness itself.

Consequently, Plotinus asserts that the even in this universe souls can be happy: highlighting, by this, an important aspect of the pagan conception of existence. If ``virtue'', as the dominating spiritual actuality, implies power, one can conceive as little that the ``good'' is accompanied by ``happiness'', as glory is separable from victory. Whoever is defeated by external or internal constraint is not ``good'': and for such a being to be happy would be unjust. But he accuses only himself of that, not the world.

Otherwise, we understand, the thing is for whoever reduces ``virtue'' to a simple ``moral'' disposition.

He certainly says, then, that ``my kingdom is not of this world'' and expects God to give happiness in the ``afterlife'' as a reward to the ``just'', who, lacking power, have tolerated and endured injustice with humility and resignation in this life. The warrior and heroic truth of ancient classical-Aryan man was different. If ``evil'' and all its materialization in violence and limits of inferior forces and corporeal things has roots in a state of degradation of the good---it is inconceivable, it is logically contradictory that it might endure as a principle of unhappiness and servitude in respect to whoever destroyed that root, having become ``good.'' If ``good'' \emph{is}, then ``evil''---suffering, passion, slavery---cannot be. They, then, will say that ``virtue'' is still imperfect; ``being'' is still incomplete; ``purity'' and ``unity'' are still ``adulterated''.

\begin{quotex}
There are those who are unarmed. But who has weapons, fights---there is no God who fights for those who are not in arms. The law requires that victory in war is to the brave, not to those who pray.

It is just that the cowardly are dominated by the wicked .
\end{quotex}


New reaffirmation of the virile, Roman, warrior spirit of the pagan tradition. New contrast with the mystical-religious attitude. New contempt for those who deprecate the ``injustice'' of earthly things and instead of blaming their cowardice, they either resign themselves to their powerlessness, or they blame the Whole or hope that ``Providence'' cares about them.

``There is no God who fights for those who are not in arms.'' This is the anti-Christian cornerstone of every warrior morality; and it carries back to the concepts explained above about the identification---from the metaphysical point of view---of ``reality'', ``spirituality'', and ``virtue''. The cowardly cannot be good, ``good'' implies a soul of a hero. And the perfection of the hero is the triumph. Ask God for victory, it would be equivalent to ask Him for ``virtue'': since victory is the body in which the very perfection of ``virtue'' takes place.

The soldiers of Fabius, while leaving, vowed not to conquer or die, but vowed to fight and to return as victors. And the victors returned. The spirit of Rome reflects the spirit of this same wisdom.

\begin{quotex}
For fear, total suppression, the soul has nothing to be afraid of.

Those who fear nothing have not reached the perfection of virtue. He is mediocrity.

In the higher man \emph{spoudaus} , feelings are not displayed as in the others. Do not reach the interior: they are other things, are suffering and grief, his or others. This would be weakness of the soul.

If suffering passes the measure---that pass it. The light that is in him will last, like the lamp of a lighthouse in the swirling wind and storm.

Master of himself even in this state, will decide what there is to do.

The \emph{spoudaus} would not be such, if a demon was acting within his action. In him, it is the sovereign mind (nous) that acts.
\end{quotex}


Plotinus admits that the superior man can sometimes have involuntary and unreflective fears, but almost as movements that do not belong to him and because his spirit is not present.

\begin{quotex}
Coming back to himself, he will drive them away ... Like a child who remains tamed only by the majesty of those who gaze on him fixedly.
\end{quotex}


Concerning suffering, it can at most cause the separation of a part of oneself not yet free from passion: but never the overturning of the higher principle. ``He will decide what there is to do.'' When it is the case, he may also withdraw from the game. Do not forget that according to Plotinus the superior man is in himself his own ``demon'' and he lives here below like an actor who plays a part he freely chose. Against the Gnostics Christians he retorted sharply, ``Why despise this world, where you yourselves have come to by your own will? It permits you to go away if good is not found there.''

As \emph{nous} in man can be precisely defined as the ``being'' principle consisting of pure intellectuality, it is the ``Olympic mind'' with respect to which the ``soul'' principle (\emph{psyche}) already represents a marginal wrapping up: at most, it is a profundity that is hidden and latent. But then, more than the ``I'', it is a ``demon'' who acts in every action. Plotinus says exactly that everything that happens without deliberation, unites a demon to a god. We now see who the opposite condition is revealed.

\begin{quotex}
There, the why of being ... it does not exist as why, but as being. Better: the two things are one.

May everyone be himself.

May our thoughts and our actions be our own. May the actions of each being belong to him. Whether they be good---whether they be evil.

When the soul has pure and impassible reason for a guide, in full self-mastery, where he wants to direct his energy. Only then can the act be said to be ours, not another's: from interiority of the soul as a purity, as a pure dominating and sovereign principle... not by action misled by ignorance and fragmented by desire ... For, then, passion, and not act, it would be in us.

Feelings are the visions of the soul asleep.

Everything of the soul that is in the body sleeps. To come out of the body, is the true awakening. Changing lives with a body, is to pass from one sleep to another sleep, from one bed to another bed.

To truly awaken, is to leave the world of bodies.
\end{quotex}


As materiality is the state of swoon of the spirit, so the reality of sleep is every reality that appears to us in the midst of the material senses. The coming out from the body and the abandonment of the world of bodies is, however, not grossly interpreted: it is essentially about an interior transformation, a self-integration in the ``intellectual nature without sleep''. And this is the true initiatic and metaphysical realization.

Very effectively, Plotinus assimilates the transformation of the body in passing from one bed to another. When it also had a consistency, the doctrine of reincarnation could not be better stigmatized, as in the part of this pagan initiate. One form is equivalent to another in the ``cycle of births'' in respect to the center which is equally distant from every point on the circumference. Metaphysical realization is a fracture in the series of conditioned states: one opening above a radically heterogeneous direction. One does not reach it by following the trail of those types who ``flee'', those who pursue a goal that they have outside themselves: in the becoming of the world of bodies.

\begin{quotex}
What must there be in front of you as a spectacle, if one looks at the outer world. But now, you need to look at yourself; to make yourself one with what you have to contemplate; to know that what you have to contemplate is yourself.

And that is your task. As someone who was invaded by the god Phoebus or a Muse, he would see himself shine in divine clarity, if he had the power at the same time to contemplate in himself this divine light.
\end{quotex}



\flrightit{Posted on 2013-03-19 by Aeneas}

\begin{center}* * *\end{center}

\begin{footnotesize}
\begin{sffamily}
\texttt{David on 2013-03-17 at 21:12 said:}

Interesting. But I always have problem with that kind of reflection as to wether it is Solar, or Promethean in a affirmation so close to existentialism and self-proclamation of God. Of course, the Self is the Whole (Atman is Brahman), and therefore everything is kept in check. But the line is thin. Still, the «Center» and «Unify yourself» themes remind me of Rumi : «You are not a drop in the ocean.

You are the entire ocean in a drop.»

\hfill

\texttt{Jason-Adam on 2013-03-17 at 21:40 said:}

This text is NOT anti-Christian, despite what moderns may think, I must point out as the goal of perfection -- of union with God -- is the goal of Catholic mystical theology. St Dionysius and the German mystics of the XIVth century are all in accord with Plotinus -- who by the way through St Augustine is an important influence on Catholic philosophy.

Though he never came out and identified as such, Baron Evola was a Neoplatonist, as some have pointed out recently. In fact Traditionalism as a whole has been described as the merger of Neoplatonism with Hinduism, and I am in accord with that interpretation.

\hfill

\texttt{David on 2013-03-17 at 21:42 said:}

What about the monistic view of Vedanta, or some mystics (sufi or christian), opposed to the dualism of body and soul of Platonic view ?

\hfill

\texttt{Logres on 2013-03-18 at 09:13 said:}

``An image of an image'' : excellent. So much for humanism today. Could someone of a lower caste, by following assiduously the rites, possibly obtain a pre-mortem initiation through obesciance?

\hfill

\texttt{scardanelli on 2013-03-18 at 09:22 said:}

I don't think that this text is inherently anti-Christian, yet after being exposed to other writers such as Tomberg and Mouravieff and coming back to Evola, it does come off as somewhat proud and vainglorious. At times it seems that Evola defines the solar path as this sense of pride, whereas in Tomberg for example, the solar path has more to do with chastity and love. I have a feeling though that this has more to do with differeing perspectives, one of priestly contemplation and the other warrior and active, rather than any dispute over metaphysics. Perhaps this is why the active/warrior path will not take one the full distance to liberation, because this vaingloriousness is essential to the active path, yet is also a major stumbling block, as described in The Charioteer in MOT... .just conjecture though. Am I right in this?

\hfill

\texttt{scardanelli on 2013-03-18 at 09:43 said:}

Of course, i'm assuming that this was written by Evola and not one of the other members of the Ur Group.

\hfill

\texttt{Cologero on 2013-03-18 at 12:48 said:}

I don't know who wrote it, as it is attributed to ``Plotino'' himself.

\textit{Introduzione alla Magia:} \url{http://books.google.com/books?id=JLOnuzXtx2UC\&printsec=frontcover\&dq=introduzione+alla+magia\&hl=en\&sa=X\&ei=M0NHUd-yDYba8wSYoIGIAQ\&ved=0CDMQ6AEwAA\#v=onepage\&q=Plotino\&f=false}

\hfill

\texttt{Cologero on 2013-03-18 at 13:00 said:}

The tone may appear over the top, but keep in mind that the ``gods'' referred to are simply the super-human states. We would refer to them, as does Rene Guenon, as angels rather than ``gods''. I don't see how that can be accomplished passively (although it may happen as a grace); rather a man must make the necessary efforts. Try to read it in that light.

I don's suspect that Plotinus is referring to the ``One'' in the same way. And to be consistent with all the other writings that have appeared here, I don't believe he is saying that man qua man is raised up to the ``gods'', but rather that the human, all too human, state must give way to the higher state.

On the other hand, and I can't remember the source of this, I believe Dostoevsky said that the man who wants to to make himself God succeeds in becoming not even an insect.

\hfill

\texttt{Saladin on 2013-03-18 at 20:19 said:}

``Though he never came out and identified as such, Baron Evola was a Neoplatonist, ... ''

I think you are right as much as we can put a label ``..ist'' on the Maestro. His praise for Emperor Julian (so called ``apostate'') who was a thorough Neoplatonist is proof of your conjecture.

\hfill

\texttt{scardanelli on 2013-03-20 at 09:56 said:}

I see... Thanks for the clarification. I think I am beginning to understand the way in which one approaches these texts. If one wants to see difference and discord, he will remain in the realm of discursive thought and will have difficulty ``penetrating'' their truths, so to speak. I suppose this is what Tomberg means when he states that the Hermetic path seeks the middle way in neutralizing binaries above.

P.S. Forgive the constant references to Tomberg and Mouravieff... they are my only reference points at this stage though i'm sure they are old hat to the rest here.

\hfill

\texttt{Jason-Adam on 2013-03-19 at 08:38 said:}

This piece needs further commentary. I am sensing a Nietzschean tone throughout the essay and Evola may be trying to remold Plotinus into a proto-Nietzsche ?

\hfill

\texttt{Cologero on 2013-03-19 at 09:17 said:}

It may be helpful for someone to verify the Plotinus quotes, since the process from Greek to Italian and then English seems to have produced some awkward wordings. I previously provided the link to the text where the references can be found. I have marked the actual quotes next to a blue vertical bar.

\hfill

\texttt{Cassiodorus on 2013-03-19 at 12:29 said:}

Do the adherents of Tradition regard the Platonic philosophical mysticism of late antiquity as a Revelation or a merely ``human system''? Does the doctrine of the Self suggest that Revelation is not, well, ``necessary''? Does the formal ``avataric descent'' actually make intellection/realization possible? That is, does it ``open the door''? Or does Revelation simply point the way?

(Apologies for beating this horse) On Plotinian emanationalism... .I'm not sure- historically speaking, traditional Christians have defended the view that \emph{creatio ex nihilo} implies that, between creator and creation, there is an absolute ontological chasm. The classical theist wouldn't, of course, deny Divine immanence, but there is an insistence on the radical divide between the natural and supernatural -- to safeguard His transcendence. To suggest anything that resembles the ``Supreme Identity'' violates creatio ex nihilo and automatically leads to howls of gnosticism and pantheism. Would it be fair to say that, from the perspective of classical theism, it is essential to distinguish between \emph{imago Dei} and the nondual doctrine of the Self?

\hfill

\texttt{Jason-Adam on 2013-03-19 at 17:58 said:}

As a Catholic who believes in the genuineness of the works of St Dionysius, I will say that despite what the Traditionalists would have us believe it was actually Neoplatonism that was influenced by Christianity -- Dionysius predating Plotinus.

Take it from where it comes from, but Dugin in his METAPHYSICS OF THE GOSPEL claims that Orthodox Christianity teaches manifestionism not creationism as the Western church does.

\hfill

\texttt{Cologero on 2013-03-19 at 17:59 said:}

From the Summa Theologica:

\begin{quotex}

\emph{creation, which is the emanation of all being}, is from the not-being which is nothing
\end{quotex}

For the whole discussion on this topic, refer to: \textit{Question 45. The mode of emanation of things from the first principle}\footnote{\url{http://www.newadvent.org/summa/1045.htm}}.

As for ``resembling'' the ``supreme identity'', by which I presume you mean God, I consider ``image and likeness'' to be synonyms, so your accusation of a violation misses the mark (allusion intended).

I hope we will not get bogged down in a theological ``debate'', which is on the level of discursive reason. There is plenty of ``real'' material to ponder, such as the meaning of being ``awake'' vs sleep, the identity of being and power, what is the ``nous'' that transcends discursive reasonging, what exactly is that ``interior transformation'' that was mentioned, and so on. I had to put up with heresy hunters from the neo-pagans, and I won't tolerate it from the other side. Gornahoor is entirely optional and I have no desire to convince anyone of anything against their will.

\hfill

\texttt{Jason-Adam on 2013-03-19 at 20:38 said:}

So creation = emanation and all things do contain a trace of the Trinity.

There is then no conflict between Plotinus and St Thomas.

\hfill

\texttt{Logres on 2013-03-19 at 23:06 said:}

It's interesting that Clement of Alexandria embraces, even claims, the term Gnostic for advancement in the Christian faith. When he debates points with them, he usually uses their level of reasoning and accepts many, if not most of their dogmas, in some form.

\hfill

\texttt{Logres on 2013-03-19 at 23:08 said:}

I think Elder Thaddeus (Your Thoughts Determine Your Life) draws a distinction between image and likeness, but only for those who are not on the true path (which Clement would say culminates in Gnosis). We all have His image, but through sin we deform His likeness. Ideally, you have both.

\hfill

\texttt{Cassiodorus on 2013-03-20 at 01:08 said:}

Colegero,

I appreciate the link that you provided.

\hfill

\texttt{Mihai on 2013-03-20 at 04:40 said:}

Trying so hard to equate Christianity and Neo-Platonism is a mistake. So is taking everything uncritically and unconditionally from the ancient greek philosophers.

For me, Plato is an invaluable resource on the path towards Truth. However, I do recognize that, every now and then in his writings, he does show inaccuracies, inconsistencies or exaggeration. It is a question of taking what is good and not being limited by some of the few failings he has. And what is good is precisely what conforms to metaphysical principles, in other words, the exposition of genuine traditional doctrine.

The same goes for Neo-Platonism. Both Plato's writings and those of Plotinus (and the rest of the neo-platonists) represent a synthesis of ancient and traditional metaphysical doctrines, but they are not direct Revelation- such as Christianity is. This means that individual opinion is sometimes present alongside with the great truths they expose, thus presenting certain limitations.

This is why they should be used as templates, as means to an end, but not ends in themselves.

The same goes with the writings of those such as Guenon, Evola and the like.

\hfill

\texttt{Jason-Adam on 2013-03-20 at 07:16 said:}

I agree, Mihai.

As I was pointing out, Plotinus was influenced by the Christian St Dionysius.

\hfill

\texttt{Logres on 2013-03-20 at 07:59 said:}

You need to spell this out more: what about natural revelation, the inner image of Man, etc.? If Christianity and the Bible are to be believed, these things have a much stronger place than traditionally accorded within the Faith, although Orthodoxy does better than most. Revelation, I agree, is more than the limiting factor of secular Reason, or a purely esoteric and private truth, however, ``Revelation'' as usually used within the Faith is exoterically denoted: Plotinus is not speaking of this, here. Additionally, there is the fact of works like the Stromata, in which Clement claims that a natural revelation actually preceded the Scripture, and is endorsed (for example) when St. Paul quotes the poet Aratus. In the Stromata, Clement goes so far as to claim that the Greek ``revelation'' is a corruption of a ``barbarian'' and ``Northern'' revelation that comes even more directly from God, and finds its natural heir in the Israelite mythos.

\hfill

\texttt{Mihai on 2013-03-21 at 06:41 said:}

Of course, the Greek myths and legends were echoes of much more ancient, perhaps even primordial revelations. There are, of course, many Hyperborean allusions in them.

However, the greek religion(s), as we know them, historically, certainly degenerated into superstitions and polytheism, with sorcery occupying top place, at least in the first centuries AD.

As for Plato and Plotinus: it is said that Plato, his main influence, traveled to Egypt to learn the ancient wisdom, which explains much of the contents of his writings. Plotinus was, of course, also influenced by ancient Eastern wisdom. Both, however, seem to have, afterwards, relied on their own efforts to put all this into a workable template. It is not clear wether or not they received some kind of initiation. In one of his books on initiation, Guenon claimed that Plotinus ``seems to have had some kind of hindrance that prevented his initiation from becoming fully active and workable'', but I don't know where he got that idea from and why.

Anyway, neither of them worked within the confines of a regular traditional revelation, in the light of which to balance their doctrines, which makes them fallible from time to time.

\hfill

\texttt{Jason-Adam on 2013-03-21 at 07:16 said:}

Plato studied under the Israelites and merged what he learned from them with his own creations.

\hfill

\texttt{Cologero on 2013-03-22 at 09:35 said:}

The point I was trying to make, Cassiodorus, even if harsh, is that matters are not so simple. As you see, Aquinas did not make a sharp distinction between creation and emanation. Here, we are taking Guenon at his word that theological propositions can be rephrased as metaphysical statements, hence we are taking that approach. If done correctly, there should be no conflict between them. To oppose one to another, would be like comparing a color to an odor; that is why we can't be bogged down in tedious and unresolvable theological disputes. As Aquinas himself put it, there are certain things he knows through metaphysics that common people must believe by faith. Other things, he leaves to faith. However, as we have often pointed out, Solovyov claims that even those dogmas can be known. That is something we will get to eventually.

The advantage of metaphysics is that more can be encompassed without unnecessary partisanship (not that it is all unnecessary).

\end{sffamily}
\end{footnotesize}

